% !TeX spellcheck = ru_RU
% !TEX root = vkr.tex

\section*{Введение}
\thispagestyle{withCompileDate}

% TODOV: проверить формулировки
% \todo{Написать про особенности: буквенные титла, особая диакритика...}
% \todo{TODOP: дописать про сложности для экспертов
% \todov{(ссылку бы на Цифр Летопис)}
% \todo{, (?направленного на создание набора отечественного ПО, / помощь в оцифровке текстов РАН? ),}

С развитием технологий компьютерного зрения и машинного обучения стало возможным во многом ускорить и автоматизировать процессы оцифровки различных текстов, хранящихся в архивах,~--- не только печатных, но и рукописных.

Однако на данный момент технологии не позволяют оцифровывать древние рукописные кириллические тексты с абсолютной точностью: периодически возникают ошибки распознавания, и без супервизии эксперта (например, славяниста или библеиста), в этом процессе никак не обойтись~--- существует необходимость редактировать печатный текст с использованием компьютерных технологий. Проблематика области заключается в том, что набор символов в старославянских и церковнославянских текстах специфичен и не поддерживается \enquote{из коробки} известными текстовыми редакторами, такими как \word{}, \librewriter{} и пр. Ввиду этого эксперту приходится тратить много времени на настройку рабочего окружения и выполнение монотонных ручных действий: поиск, копирование и дальнейшая вставка нужного кириллического символа; вычитка и исправление \enquote{сломавшихся} символов при смене шрифта.

В рамках проекта \enquote{Цифровой летописец} совместно отделом рукописей Российской академии наук производится оцифровка нескольких кириллических рукописных текстов, в т.ч. книги Сираха \cite{sirahPhotoYaDisk}, и было принято решение разработать инструмент, позволяющий упростить набор и редактирование оцифрованных кириллических текстов.

% \todo{Добавить проблематику: нет приложения для удобной работы с оцифрованными кириллическими текстами, сложности представления кириллических рукописей в виде текста \unicode{}}

% \todo{Добавить про то, что будет представлено в данной работе: сравнительный обзор шрифтов, приложение для работы с кириллическими рукописями}
